% ============================================================
%  v3_compliance.tex — Referee Compliance (Section 5)
%  V3 draft: How the directive propagated through the referee workforce
%  Numbers sourced from a06 analysis + results_v3_compliance.json
% ============================================================

\section{Referee Compliance}\label{sec:compliance}

The first-stage results in Section~\ref{sec:firststage} establish that
the directive raised average stoppage time. But aggregate effects can
mask heterogeneous implementation: if compliance is concentrated among
a few zealous referees while others ignore the directive, the policy
implications differ sharply from a setting of near-universal adoption.
This section examines how the directive propagated through the referee
workforce, drawing on 25,464 referee-match observations covering 407
unique referees.

\subsection{Referee-Level Compliance}

We measure individual compliance as the share of treated-league
referees whose average post-directive stoppage time exceeds their own
pre-directive average. At the most conservative threshold---restricting
to referees with at least 20 matches in both periods---100\% of
85 referees increased their stoppage time. Relaxing the threshold to 10
matches yields 98\% (100 of 102), and the unrestricted sample (all referees with
both pre- and post-directive observations) yields 95\% (114 of 120). The mean
within-referee increase is 1.46 minutes.

These compliance rates are remarkable by the standards of the
regulatory compliance literature. \citet{cabana1999compliance} surveys
physician adherence to clinical practice guidelines and finds rates
between 20--100\%, with a median around 50\%.
\citet{cohen1990compliance} documents similarly mixed results for
environmental regulation. Even \citet{duflo2013enforcement}, studying
third-party environmental audits in Gujarat, India---a setting with
formal monitoring and financial incentives---finds compliance rates
well below 90\% in the initial periods. By contrast, the FIFA
directive achieved near-universal adoption through what appears to be
focal-point coordination: the directive shifted the common expectation
of ``how much stoppage time is appropriate'' without specifying a
numerical target.

A heterogeneity analysis by subgroup confirms that compliance was
universal in direction and broadly similar in magnitude.
Figure~\ref{fig:compliance_forest} reports within-referee treatment
effects estimated separately by league, experience tercile, and
pre-directive stoppage-time baseline. All subgroups show large,
statistically significant increases ($p < 0.001$ in every case).

Experience does not produce the gradient one might expect.
Low-experience referees show the largest shift ($+1.56$~minutes,
SE~$= 0.073$), while medium- ($+1.28$, SE~$= 0.081$) and
high-experience ($+1.30$, SE~$= 0.120$) referees show somewhat
smaller adjustments. This pattern is consistent with experienced
referees already operating closer to the new norm pre-directive, leaving
less room for upward correction. Pre-directive baseline style tells a
similar story: ``low baseline'' referees adjusted the most ($+1.65$~min),
while ``high baseline'' referees adjusted the least ($+1.50$~min),
confirming that the directive primarily moved the lower tail of the
distribution upward.

\subsection{Match-Level Compliance}

Individual compliance rates measure whether referees moved in the
right direction. A stricter test asks whether the post-directive
stoppage time in each \emph{match} exceeds the pre-directive league
median---that is, whether the directive shifted the entire
distribution, not just individual means. By this measure, 70.2\% of
post-directive matches in treated leagues exceed the pre-directive
league-specific median stoppage time. Under the null hypothesis of no
effect, the expected exceedance rate is 50\%, so the observed
rate represents a 20~percentage-point excess. This match-level measure
is more conservative than the referee-level rate because it is
sensitive to within-referee variance: a referee who increased on
average but has high match-to-match variability will have some matches
below the pre-directive median.

\subsection{Convergence: Compression of the Referee Distribution}

If compliance were heterogeneous---with some referees adding five
minutes and others adding none---we would expect between-referee
dispersion to increase post-directive relative to the new mean.
Figure~\ref{fig:compliance_convergence} shows that the between-referee
coefficient of variation (CV) of average stoppage time \emph{fell}
from 0.173 to 0.144 after the directive, even as the raw standard
deviation rose modestly (0.87 to 0.93~min) due to the higher mean
(5.06 to 6.47~min). Referees who had been stingy made large upward
corrections, while those already generous adjusted by smaller
increments. Relative to the scale of post-directive stoppage, the
distribution compressed, producing convergence rather than divergence.

This convergence pattern has a natural interpretation in the
regulatory compliance framework of \citet{duflo2013enforcement}. The
FIFA directive operated in a high-monitoring environment: every
top-flight match is televised, referees' decisions are reviewed by
assessors, and post-match data are publicly available. In such
environments, non-compliance is immediately observable, making formal
sanctions largely unnecessary. The directive functioned as a
coordination device, shifting the focal point of acceptable referee
behaviour. Referees who had previously added minimal stoppage time now
faced a clear signal that the norm had changed, and the public nature
of their decisions made deviation costly in reputational terms.

\begin{figure}[htbp]
  \centering
  \includegraphics[width=0.85\textwidth]{fig_compliance_convergence_v3}
  \caption{Convergence of Referee Stoppage-Time Distributions}
  \label{fig:compliance_convergence}
  \begin{minipage}{0.85\textwidth}
  \footnotesize
  \textit{Notes:} Between-referee coefficient of variation (CV) of mean
  second-half stoppage time, computed at the referee-season level.
  The vertical dashed line marks the directive onset (2022--23).
  Treated leagues are shown in dark grey; control leagues in light grey.
  The decline in CV for treated leagues indicates relative convergence
  among referees in response to the directive, even as the raw standard
  deviation rose modestly with the higher post-directive mean.
  \end{minipage}
\end{figure}

\subsection{Adaptation Speed}

A natural question is whether compliance emerged gradually or
arrived as an immediate step change. We decompose the within-referee
treatment effect by post-directive season. The first full
post-directive season (2023--24) shows an increase of $+1.26$~minutes
in treated-league referees relative to their pre-directive average.
The second season (2024--25) shows $+1.31$~minutes, and the third
(2025--26) shows $+1.38$~minutes. The flat profile rules out both
gradual ramp-up and novelty-driven decay: referees adopted the new
norm immediately and sustained it without erosion. This immediate-jump
pattern is consistent with the focal-point interpretation---the
directive did not require referees to learn a new skill, but rather to
recalibrate a single parameter (how much time to add) in a direction
most already understood.

\begin{figure}[htbp]
  \centering
  \includegraphics[width=0.85\textwidth]{fig_compliance_forest_v3}
  \caption{Referee Compliance by Subgroup}
  \label{fig:compliance_forest}
  \begin{minipage}{0.85\textwidth}
  \footnotesize
  \textit{Notes:} Forest plot showing compliance rates and within-referee
  DiD estimates across subgroups. Left panel: share of referees whose
  post-directive average exceeds their pre-directive average.
  Right panel: within-referee DiD coefficient with 95\% CI. Subgroups
  defined by experience tercile, league, and pre-directive stoppage
  quartile.
  \end{minipage}
\end{figure}

\subsection{Discussion}

The compliance results reinforce the first-stage findings through a
different lens. The aggregate DiD coefficient ($+0.734$, Table~\ref{tab:first_stage})
reflects not a few outlier referees but a near-universal shift in
enforcement behaviour. Three features of the compliance pattern are
noteworthy for the broader regulatory compliance literature.

First, the 95--100\% compliance rate was achieved without formal
sanctions or monitoring beyond what already existed. The directive was
a memo, not a rule change with associated penalties. This aligns with
models of norm-based compliance in which the principal need not
monitor every agent---the common knowledge that the norm has shifted
is sufficient \citep{johnson2020compliance}.

Second, the convergence in referee behaviour suggests that the
pre-directive distribution reflected heterogeneous interpretations of
an ambiguous norm (``add time for stoppages'') rather than genuine
disagreement about optimal match length. When the norm was clarified,
referees converged, implying that much of the pre-directive variance
was coordination failure rather than preference heterogeneity.

Third, the immediate adoption and absence of decay over three seasons
suggest that the new norm has become self-sustaining. Once all
referees add more time, any individual referee who reverts to the old
norm would stand out. The equilibrium has shifted, and the directive
has become---in regulatory parlance---a \emph{self-enforcing}
standard.
